\chapter{Models, Messages and Runnables}
\label{ch:models}

\begin{quote}\itshape
One interface, five methods, and an \texttt{a} prefix. That is nearly the whole
of it.
\end{quote}

\noindent
Underneath every component in this ecosystem is one interface with a small,
regular surface. Learn it once and you can drive a model, a tool, a graph or an
entire agent without reading their documentation --- which is a larger claim than
it sounds, and it is why Chapter~\ref{ch:landscape} could show you that an agent
is a \api{Runnable} and have that mean something.

This chapter also covers the part of 1.0 that quietly deleted the most code from
real projects, and which almost nobody lists among the headline features:
standard content blocks.

% ============================================================
\section{The Runnable contract}
\label{sec:runnable}
\index{Runnable@\texttt{Runnable}}
\index{Runnable@\texttt{Runnable}!batch}
\index{asyncio!a-prefix convention}

\mermaidfig{lc-runnable-surface}{%
  The whole interface. Everything in the ecosystem implements it.}{%
  The Runnable surface and the async naming rule.}

\begin{center}
\small
\begin{tabularx}{\linewidth}{@{}llX@{}}
\toprule
\textbf{Sync} & \textbf{Async} & \textbf{Does} \\
\midrule
\code{invoke} & \code{ainvoke} & One input, one output \\
\code{batch} & \code{abatch} & Many inputs, with internal concurrency \\
\code{stream} & \code{astream} & Chunks as they are produced \\
\code{stream\_events} & \code{astream\_events} & Lifecycle events, and
  Chapter~\ref{ch:interrupts}'s v3 projections \\
\bottomrule
\end{tabularx}
\end{center}

\textbf{The naming rule: the async method is the same name with an \code{a} in
front.} It holds across the whole ecosystem --- \code{ainvoke},
\code{astream}, \code{aget}, \code{aput}, \code{abefore\_model}. Once you have
noticed it you stop looking up async variants.

There are also wrappers that take a \api{Runnable} and return another one:

\begin{itemize}[leftmargin=1.4em,nosep]
\item \code{with\_fallbacks} --- try a different component on failure
\item \code{with\_retry} --- try the same one again
\item \code{with\_config} --- pin configuration
\item \code{bind} --- fix an argument
\end{itemize}

\medskip
\noindent
That closure property is what makes composition work at all.

\begin{dotnetbox}
The closest analogue is \code{IEnumerable<T>} and LINQ: one small interface that
everything implements, plus operators that take one and return another. The
payoff is the same --- you learn the shape once rather than learning each
component separately.
\end{dotnetbox}

\begin{warning}
\code{batch} runs its inputs concurrently \emph{internally}, which is convenient
and is also an unbounded fan-out at your provider. Chapter~\ref{ch:asyncio}
measured what a cap costs and Chapter~\ref{ch:serving} explains why it must be
sized per process. Use \code{max\_concurrency} in the config, or do not use
\code{batch} for anything large.
\end{warning}

% ============================================================
\section{Choosing a model}
\label{sec:init-chat-model}
\index{init\_chat\_model@\texttt{init\_chat\_model}}
\index{model!provider selection}
\index{packages!provider integrations}

\begin{python}[caption={Two ways to name a model}, label={lst:init-model}]
from langchain.chat_models import init_chat_model

model = init_chat_model("anthropic:claude-sonnet-4-5", temperature=0)

# or directly, when you need provider-specific arguments
from langchain_anthropic import ChatAnthropic
model = ChatAnthropic(model="claude-sonnet-4-5", max_tokens=4096)
\end{python}

The \code{provider:model} string form is the one to prefer in application code,
because the provider becomes configuration rather than an import. Each provider
lives in its own distribution --- \pkg{langchain-openai},
\pkg{langchain-anthropic} and so on --- which is why
Chapter~\ref{ch:landscape}'s package map has a whole category for them.

\begin{note}
Provider-neutral is a real property of this ecosystem and it is worth preserving
deliberately. Every listing in this book works against any provider with tool
calling. The mini-project reads its provider from configuration and nothing in
\code{src/} imports a provider package directly.
\end{note}

% ============================================================
\section{Messages, and the thing 1.0 fixed}
\label{sec:content-blocks}
\index{content\_blocks@\texttt{content\_blocks}}
\index{message!types}
\index{message!content blocks}

Five message types --- system, human, AI, tool, and function (legacy). The
interesting part is what an \api{AIMessage} carries:

\begin{python}[caption={What comes back}, label={lst:message-surface}]
response.text              # the text, as a string
response.content_blocks    # normalised blocks -- the 1.0 addition
response.tool_calls        # what the model wants to run
response.usage_metadata    # input/output tokens -- the whole cost story
\end{python}

\begin{console}
>>> AIMessage(content="hello").content_blocks
[{'type': 'text', 'text': 'hello'}]
\end{console}

\mermaidfig{lc-content-blocks}{%
  What \api{content\_blocks} removed: the provider-shaped branch in your
  application.}{%
  Provider-specific responses normalised into content blocks.}

Before 1.0, handling reasoning output from one provider and citations from
another meant provider-specific branching in application code, because
\code{content} was a string for some providers and a list of dictionaries for
others. \api{content\_blocks} normalises it: \code{text}, \code{reasoning},
\code{citation}, \code{tool\_call} and the multimodal types, with the same shape
regardless of who produced them.

\begin{note}
This is the underrated headline of 1.0. It is not a new capability --- it is the
deletion of a category of code. If you are migrating and you have a module named
something like \code{response\_parsing.py} with a branch per provider, it exists
to solve a problem that no longer exists.
\end{note}

\subsection{Counting what it cost}
\label{sec:usage-metadata}
\index{cost!usage metadata}
\index{usage\_metadata@\texttt{usage\_metadata}}

\code{usage\_metadata} carries input and output token counts per call. It is the
foundation of everything in Chapter~\ref{ch:context} and
Chapter~\ref{ch:security}, and it costs nothing to record.

\begin{warning}
Put cost accounting in from the first model call, not later. Chapter~\ref{ch:context}
measured a 5$\times$ difference between context strategies --- a difference you
cannot see, let alone act on, without per-call token counts already being
recorded. Retrofitting them means re-running everything you wanted to compare.
\end{warning}

% ============================================================
\section{Reliability at the model boundary}
\label{sec:fallbacks}
\index{with\_fallbacks@\texttt{with\_fallbacks}}
\index{reliability!fallbacks}
\index{retry}

\begin{python}[caption={Degrading instead of failing}, label={lst:fallbacks}]
robust = primary.with_fallbacks([secondary, tertiary])
\end{python}

Fallbacks fire on exception: if the primary raises, the next is tried with the
same input. \code{with\_retry} retries the \emph{same} model, which is the right
response to a transient error and the wrong response to an outage.

\begin{center}
\small
\begin{tabularx}{\linewidth}{@{}lX@{}}
\toprule
\textbf{Failure} & \textbf{Answer} \\
\midrule
Transient error, rate limit & \code{with\_retry}, or the provider client's own
  retry --- with jitter (Chapter~\ref{ch:serving}) \\[3pt]
Provider outage & \code{with\_fallbacks} to a different provider \\[3pt]
Context overflow & Chapter~\ref{ch:context}, not a retry \\[3pt]
The model answered badly & Evaluation (Chapter~\ref{ch:observability}), not a
  retry \\
\bottomrule
\end{tabularx}
\end{center}

Chapter~\ref{ch:agents}'s model-fallback middleware does the same thing inside an
agent loop, which is where you usually want it.

\begin{warning}
A fallback to a cheaper or weaker model silently changes your quality
distribution, and it does so exactly when you are least able to notice --- during
an incident. Record which model actually answered, and make it a dimension in
your evaluation set (Chapter~\ref{ch:observability}).
\end{warning}

% ============================================================
\section{Streaming, multimodality and reasoning}
\label{sec:streaming-basics}
\index{streaming!tokens}
\index{multimodality}
\index{reasoning content}

\begin{python}[caption={The smallest streaming loop}, label={lst:astream}]
async for chunk in model.astream(messages):
    print(chunk.text, end="", flush=True)
\end{python}

That is the primitive; Chapter~\ref{ch:interrupts} builds the agent-level
projections on top of it, and Chapter~\ref{ch:serving} puts it behind HTTP.

\textbf{Multimodality} arrives through the same content blocks: an image or a
document is a block in the message rather than a separate API. Which providers
accept which block types varies, and that is one of the few places where the
abstraction genuinely leaks --- Appendix~\ref{app:packages} has the matrix.

\textbf{Reasoning models} return their thinking as a \code{reasoning} block. Two
things follow that are easy to miss: reasoning tokens are \emph{billed}, and they
are usually not what you want to show a user. Read the block type rather than
concatenating everything in \code{content}.

\begin{versionbox}
Reasoning-model surfaces are the fastest-moving part of the provider integrations
--- budgets, block names and billing treatment have all changed within the 1.x
line. Check your provider package's release notes rather than trusting a
six-month-old example, including this one.
\end{versionbox}

% ============================================================
\section{What to take away}
\label{sec:ch5-summary}

\begin{itemize}[leftmargin=1.4em]
\item One interface; the async twin is the same name with an \code{a} in front.
\item Wrappers return \api{Runnable}s, which is what makes composition work.
\item \code{batch} is an unbounded fan-out unless you cap it.
\item \api{content\_blocks} normalises provider differences. If you have a
      per-provider parsing module, delete it.
\item Record \code{usage\_metadata} from the first call. You cannot retrofit a
      cost comparison.
\item \code{with\_retry} for transient errors, \code{with\_fallbacks} for
      outages. Record which model answered.
\end{itemize}

\begin{exercisebox}
\textbf{1.} Send the same prompt to two providers and print
\code{content\_blocks} from each. Confirm one code path handles both.

\textbf{2.} Build a token-and-cost accumulator over \code{usage\_metadata} and
run a five-turn conversation through it. Compare the total against your
intuition; most people underestimate, because of
Chapter~\ref{ch:context}'s re-sending.

\textbf{3.} Wrap a model in \code{with\_fallbacks} to a deliberately broken one,
then reverse the order. Confirm you can tell from the output which answered.
\end{exercisebox}

\begin{projectbox}
\textbf{Stage 02 --- streaming CLI chat.} The first real model call: a
command-line chat on \code{astream}, printing tokens as they arrive and
reporting tokens and cost per turn.

The contract in \code{stages/02-streaming-chat/README.md} asks that streamed
chunks assemble into the same message the non-streamed call returns, tested
against a fake --- and that nothing reads \code{content} as a string when
\code{content\_blocks} is available.
\end{projectbox}
